
\documentclass[12pt]{article}
%\renewcommand{\baselinestretch}{1.5}
\usepackage{sectsty,setspace}
\usepackage[top=1.87cm, bottom=1.87cm, left=1.87cm, right=1.87cm]{geometry} 
\usepackage{epstopdf} % for pdf creation
\usepackage{amsmath,latexsym,amssymb,wasysym} % improving structure
\usepackage{natbib}
%\usepackage{hyperref} % to get the little green boxes around the refs
\usepackage{times} % setting font to times
\usepackage{fancyhdr} % for custom headers/footers

\setstretch{1}
%\setlength{\baselineskip}{0.167in} 
\setlength\parindent{0pt} % no indents throughout

%%%%%%%%%%%%%%%%%%%%%%%%%%%%%%
% NSERC's formatting
%%%%%%%%%%%%%%%%%%%%%%%%%%%%%%
% Setup header
\begin{document}
\pagestyle{fancy}
\fancyhf{} % clear header/footer
\fancyhead[R]{Christophe Rouleau-Desrochers} % name on right side
\fancyfoot[C]{\thepage}
\renewcommand{\headrulewidth}{0pt}


%%%%%%%%%%%%%%%%%%%%%%%%%%%%%%

%%%%%%%%%%%%%%%%%%%%%%%%%%%%%%
% Regular formatting
%%%%%%%%%%%%%%%%%%%%%%%%%%%%%%
%\title{Improving forest carbon models under climate change using drones and experiments}
%\date{\today}
%\author{Christophe Rouleau-Desrochers}
%\begin{document}
%%%%%%%%%%%%%%%%%%%%%%%%%%%%%%

\section*{Improving Forest Carbon Models under Climate Change Using Drones and Experiments}

\textbf{Contexte :} Les changements climatiques transforment les systèmes naturels à l'échelle mondiale \cite{parmesan_poleward_1999,rosenzweig_attributing_2008}, et l'impact biologique le plus observé est changements phénologiques, soit l'étude des événements périodiques et récurents du vivant  \cite{parmesan_globally_2003,cleland_shifting_2007,lieth_phenology_1974,woolway_phenological_2021,menzel_european_2006}. Les évènements phénologiques printaniers (débourrement, floraison) avancent de 0,5 \cite{wolfe_climate_2005} à 4,2 jours par décennie \cite{chmielewski_response_2001,fu_recent_2014} et sont  l'effet de la température \cite{chuine_why_2010,cleland_shifting_2007,penuelas_responses_2001} ; les événements automnaux (formation des bourgeons, coloration des feuilles) sont retardés, mais bien plus faiblement \cite{gallinat_autumn_2015,jeong_macroscale_2014}, sous l'effet du raccourcissement de la photopériode \cite{cooke_dynamic_2012,flynn_temperature_2018,korner_phenology_2010} et du refroidissement \cite{cooke_dynamic_2012,delpierre_temperate_2016}. Ces décalages allongent la saison de croissance, et l'on admet depuis longtemps que la croissance des arbres devrait augmenter en conséquence. Pourtant, un nombre croissant de travaux ne parvient pas à établir ce lien \cite{dow_warm_2022,silvestro_longer_2023,matula_shifts_2023,charletdesauvage_temperature_2022,cabon_crossbiome_2022} : Dow \textit{et al}. (2022) ont montré que malgré un démarrage plus hâtif, ni le taux de croissance ni l'accroissement annuel n'augmentaient avec la longueur de la saison. Comme les forêts tempérées représentent 40 \% du puits net de carbone mondial \cite{gibbs_revised_2025}, ce paradoxe se répercute directement sur les projections du cycle du carbone et leurs rétroactions climatiques \cite{richardson_climate_2013,swidrak_comparing_2013}.

Résoudre ce paradoxe exige de comprendre pourquoi les arbres ne croissent pas davantage lorsque la saison s'allonge. Je pose deux classes de limites non exclusives : externes (environnementales) \cite{kolar_response_2016} et internes (physiologiques) \cite{zohner_effect_2023}. Les facteurs externes, sécheresse, gel printanier et vagues de chaleur, surviennent simultanément et agissent sur des communautés entières plutôt que sur des individus \cite{choat_triggers_2018,li_widespread_2023,intergovernmentalpanelonclimatechange_detection_2014,reinmann_compensatory_2023} ; seule l'expérimentation permet de les dissocier, par exemple un printemps chaud d'une sécheresse plus tardive \cite{morin_changes_2010,primack_observations_2015}, et d'affiner les projections de séquestration du carbone \cite{green_limits_2022,cabon_crossbiome_2022}. Les limites internes tiennent à la répartition du carbone entre croissance et réserves : si la croissance est limitée par le puits plutôt que par la source, une saison plus longue produit davantage de photosynthétats sans produire plus de bois ; le surplus s'accumule en glucides non structuraux \cite{zohner_effect_2023,chapin_ecology_1990}, et le bénéfice ne se manifeste, le cas échéant, qu'une année plus tard \cite{landhausser_partitioning_2012}. Vérifier cette hypothèse exige des données pluriannuelles, que mon expérience achevée en 2025 fournit désormais. Les expériences portent par ailleurs surtout sur de jeunes arbres, dont les réponses éclairent la régénération mais se transposent mal aux arbres matures, qui détiennent l'essentiel de la biomasse forestière \cite{augspurger_differences_2003,silvestro_longer_2023,vitasse_ontogenic_2013}. Combler cet écart in situ est maintenant réalisable : l'imagerie multispectrale par drone, couplée à l'apprentissage automatique, donne accès à des échantillons inatteignables au sol, à une résolution spatiale et temporelle plus fine que celle des satellites \cite{berra_assessing_2019,piao_plant_2019,teng_bringing_2025}. Je propose donc deux expériences testant les limites internes (chapitre 1) et externes (chapitre 2), ainsi qu'un projet d'observation à grande échelle sur des canopées matures (chapitre 3). L'originalité du programme tient à ce qu'aucune étude n'a encore testé ces deux classes de limites dans un même dispositif, ni relié les réponses des juvéniles à celles des arbres matures.

%<><><><><><><><><><><><><><><><><><><><>
% CHAPITRE 1 %
%<><><><><><><><><><><><><><><><><><><><>
\textbf{Chapitre 1. Effets différés de l'allongement de la saison, limites internes (année 1) :}
J'ai vérifié dans mon mémoire que les décalages phénologiques affectent la croissance de la saison en cours, mais aucune expérience n'a quantifié leurs effets différés sur les années suivantes \cite{chapin_ecology_1990,landhausser_partitioning_2012,schott_premature_2013}. \textit{Objectif :} déterminer si la réponse de la croissance à la longueur de la saison dépend des conditions de l'année précédente. \textit{Hypothèse :} les arbres qui prolongent leur activité une année puisent dans leurs réserves, ce qui annule les gains l'année suivante, l'effet étant le plus marqué chez les espèces qui investissent dans le stockage. À l'automne 2025, j'ai terminé une expérience sur la longueur de saison avec 15 réplicats de sept espèces suivis sur deux années consécutives. En année 1, j'analyserai les dizaines de milliers d'observations recueillies en 2024 et 2025 par un modèle bayésien hiérarchique conjoint estimant simultanément phénologie et croissance. Le modèle mettra l'information en commun entre espèces et propagera l'incertitude sur les dates phénologiques directement dans la composante de croissance ; les études actuelles traitent ces dates comme des covariables fixes, une simplification qui biaise la taille des effets. Je le développerai avec le Dr Charles Margossian, spécialiste de l'inférence bayésienne, sur la base du flux de travail Stan établi durant ma maîtrise. Les données étant déjà acquises, ce chapitre ne comporte aucun risque de collecte et livrera l'outil de modélisation réclamé par Wolkovich \textit{et al}. (2025).

%<><><><><><><><><><><><><><><><><><><><>
% CHAPITRE 2 %
%<><><><><><><><><><><><><><><><><><><><>
\textbf{Chapitre 2. Expérience de sécheresse et de gel printanier, limites externes (années 2 et 3) :}
Les changements climatiques modifieront aussi le moment des déficits hydriques et la fréquence des gels tardifs \cite{dox_wood_2022}. La dendrochronologie montre que sécheresses et gels causent d'importantes pertes de tissus \cite{kramer_why_2012,baumgarten_no_2023,dandrea_winters_2019}, mais on ignore si les arbres touchés croissent aussi moins, et si le moment du stress importe davantage que son intensité \cite{chamberlain_late_2021,baumgarten_no_2023,lian_summer_2020,zhang_drought_2021}. \textit{Objectif :} quantifier comment le moment de la sécheresse et du gel modifie la croissance chez des espèces aux stratégies d'histoire de vie contrastées. \textit{Hypothèse :} la perte de croissance augmente avec la proximité du stress par rapport au pic de croissance de l'espèce, de sorte qu'une sécheresse hâtive est largement compensée alors qu'une sécheresse au solstice ne l'est pas.
Dès l'année 2, j'utiliserai 15 réplicats de 16 espèces décidues nord-américaines (8 paires congénériques, pour éviter les effets confondants d'une histoire évolutive partagée) couvrant le gradient acquisitif--conservateur, en cinq traitements et un témoin, soit 1440 individus, un effectif comparable à l'expérience menée à terme durant ma maîtrise. Pour la sécheresse, les arbres seront placés en chambres chaudes et sèches jusqu'au point de flétrissement propre à chaque espèce, mesuré par des sondes d'humidité dans chaque pot, puis ramenés en conditions ambiantes sous irrigation constante. Les trois sécheresses ne diffèrent que par leur moment : juste après la feuillaison ; une semaine avant le solstice, période de croissance maximale chez plusieurs espèces \cite{anderson-teixeira_carbon_2021,dorangeville_drought_2018,mcmahon_general_2015} ; et en fin de saison, avant l'arrêt de croissance. Pour le gel, les arbres seront forcés à débourrer en chambres chaudes, puis maintenus à $-4^{\circ}$C une heure, au débourrement pour le premier traitement et une fois les feuilles déployées pour le second, suivant Chamberlain \textit{et al}. (2021). Dans tous les traitements, je suivrai la phénologie deux fois par semaine, estimerai la biomasse en début et en fin de saison par équations allométriques validées sur un sous-échantillon récolté destructivement, et équiperai un sous-ensemble d'arbres de dendromètres magnétiques pour résoudre le moment exact des changements de croissance. Les analyses reprendront le cadre hiérarchique bâti au chapitre 1.

%<><><><><><><><><><><><><><><><><><><><>
% CHAPITRE 3 %
%<><><><><><><><><><><><><><><><><><><><>
\textbf{Chapitre 3. Synchronie croissance--phénologie par imagerie de drone (années 1 à 3) :}
La façon dont la croissance suit la phénologie foliaire, et dont cette synchronie varie entre espèces, est déterminante pour la comptabilité du carbone forestier, mais demeure peu étudiée \cite{klein_coordination_2016,kramer_importance_2000}. \textit{Objectif :} quantifier la relation entre phénologie foliaire et saisonnalité de la croissance pour des arbres de canopée individuels en forêt mixte. \textit{Hypothèse :} les espèces diffèrent systématiquement dans ce couplage le long d'un gradient d'histoire de vie, ce qui fait de la phénologie de canopée un indicateur biaisé de la croissance. Pendant les trois années, je survolerai en drone multispectral une forêt mixte de la Station de biologie des Laurentides (Saint-Hippolyte, QC), avec des méthodes d'apprentissage automatique de pointe \cite{ball_accurate_2023,ulku_deep_2022}. Ce site me permet de poursuivre des travaux déjà menés par mon laboratoire \cite{flynn_temperature_2018} et d'établir un partenariat avec le Dr Étienne Laliberté, dont le laboratoire y travaille \cite{cloutier_influence_2024}. Je volerai tous les trois à cinq jours durant les transitions printanière et automnale, et une fois par semaine en milieu de saison, et j'utiliserai BalSAM, un modèle qui segmente les couronnes avec précision à partir d'images répétées \cite{teng_bringing_2025}, pour suivre chaque couronne d'un vol à l'autre et inférer le début et la fin de saison par espèce et par individu \cite{berra_assessing_2019,fawcett_monitoring_2021}. J'y joindrai 50 dendromètres enregistrant la variation du diamètre des tiges de canopée. Cette combinaison me permettra de tester la relation phénologie--croissance à l'échelle de l'individu, améliorant les données dont dépendent les modèles du cycle du carbone \cite{richardson_climate_2013,swidrak_comparing_2013} et les méthodes de télédétection qui les produisent.

%<><><><><><><><><><><><><><><><><><><><>
% DIFFUSION %
%<><><><><><><><><><><><><><><><><><><><>
\textbf{Diffusion et mobilisation des connaissances :}
Des saisons plus longues pourraient favoriser certaines espèces et en désavantager d'autres, remodelant la dynamique forestière nord-américaine. Je publierai les trois chapitres en libre accès et les présenterai à des congrès nationaux et internationaux. Le code, les données et les protocoles de vol et de segmentation seront archivés sur le Knowledge Network for Biocomplexity, de façon réutilisable par d'autres équipes. Je continuerai de contribuer au programme de science participative Tree Spotters, auquel je collabore depuis 2024, et démarrerai un suivi phénologique à Montréal en recrutant dans des quartiers comme Montréal-Nord et Saint-Michel, où le couvert arboré est le plus faible. Fort de mon expérience en vulgarisation, je poursuivrai la diffusion grand public en français et remettrai les résultats aux participants.
%<><><><><><><><><><><><><><><><><><><><> 
% REFERENCES
%<><><><><><><><><><><><><><><><><><><><>
\newpage

\bibliography{bibproposed.bib}
\bibliographystyle{nature} % set citation style 

\end{document}
